% Day 2 paper-and-pencil worksheet: discover the algorithm of a hypothesis test.
% Students are told what the binomial distribution is. The student pages name only
% the null hypothesis (Q2), the test statistic, and its sampling distribution under the
% null hypothesis (Q3); significance, rejection region, and above all the p-value are not
% given. Six questions on two pages: students build a decision rule (Q4), apply it to the
% data (Q5), and write the whole procedure as a recipe with step 1 pre-filled (Q6).
% The p-value is discovered at the board after Q6 (the dealer's protest; see the key's notes), together with why a tail
% rather than a single outcome, and what the number does not mean.
% The instructor key maps each question to its standard name for the lecture that follows.
%
% One source, two builds (see Makefile):
%   make student -> coin_test.pdf      (blank answer boxes)
%   make key     -> coin_test_key.pdf  (answers in the same boxes, plus an instructor-notes page)
% The key is switched on only by the build (\showsolutions); never edit the switch by hand.
% Modernized from old_ref/worksheet_NHT_discovery_v2.pdf (2025 course),
% with two rounds of pedagogical review from Codex (2026-10-05).

\documentclass[11pt]{article}
\usepackage[lecture,noheaders,nobib]{catniplab}
\usepackage{amsmath}
\input{math_preamble}
\usepackage{pgfplots}
\pgfplotsset{compat=1.18}

\newif\ifsolutions
\ifdefined\showsolutions \solutionstrue \fi

% Answer box: same height in both builds, so the key lines up page for page.
% #1 = box height, #2 = the answer, shown only in the key.
\newcommand{\answerbox}[2][2.2cm]{%
  \par\smallskip\noindent
  \begin{tcolorbox}[colback=white, colframe=gray!55, boxrule=0.4pt, sharp corners,
      height=#1, valign=top, left=5pt, right=5pt, top=3pt, bottom=3pt]
    \ifsolutions{\small\color{darkblue}#2}\fi
  \end{tcolorbox}\par}

% Answer box that starts pre-filled on both builds (a worked first line), then the answer in the key.
% #1 = box height, #2 = the pre-filled text, #3 = the rest of the answer, shown only in the key.
\newcommand{\prefilledbox}[3][2.2cm]{%
  \par\smallskip\noindent
  \begin{tcolorbox}[colback=white, colframe=gray!55, boxrule=0.4pt, sharp corners,
      height=#1, valign=top, left=5pt, right=5pt, top=3pt, bottom=3pt]
    {\small #2}\ifsolutions\par{\small\color{darkblue}#3}\fi
  \end{tcolorbox}\par}

\newcommand{\question}[1]{\par\medskip\noindent\textbf{#1}\enspace}

\setlength{\parindent}{0pt}
\setlength{\parskip}{3pt}
\pagestyle{plain}

\begin{document}

% ------------------------------------------------------------------ page 1
{\small INDP/ICDP Data Analysis, Statistics, and Experimental Design \hfill Day 2}

\medskip
{\LARGE\bfseries Is the dealer's coin fair?}\par
\smallskip
{\large Invent a way to decide}
\ifsolutions\hfill{\large\bfseries\color{c:time}Instructor key}\fi

\medskip
Name: \rule{0.5\textwidth}{0.4pt}

\medskip
Work on your own, with a pencil; about 30 minutes.
No calculator needed: every number you need is in the table, and you only add.

\begin{highlightbox}
\textbf{The binomial distribution.}
Flip a coin $n$ times, independently, with probability $p$ of heads on every flip,
and count the heads, $X$.
Then $X$ has the \emph{binomial distribution}:
\[
  P(X = k) = \binom{n}{k}\, p^k (1-p)^{n-k}, \qquad k = 0, 1, \dots, n,
\]
where $\binom{n}{k}$, read ``$n$ choose $k$'', is the number of ways to choose
which $k$ of the $n$ flips are the heads,
that is, the number of sequences of $n$ flips with exactly $k$ heads:
\[
  \binom{n}{k} = \frac{n!}{k!\,(n-k)!},
  \qquad\text{for example}\quad
  \binom{10}{2} = \frac{10 \times 9}{2 \times 1} = 45.
\]
For a fair coin ($p = 1/2$) and $n = 10$, all $2^{10} = 1024$ sequences are equally likely,
so $P(X = k) = \binom{10}{k} / 1024$.
The chart and table below give these probabilities, rounded to three decimals.
\end{highlightbox}

\medskip
\noindent
\begin{minipage}[c]{0.55\textwidth}
\begin{tikzpicture}
  \begin{axis}[
      ybar, bar width=9pt,
      width=0.93\linewidth, height=5.2cm,
      ymin=0, ymax=0.3, ytick={0,0.1,0.2,0.3}, yticklabels={0, 0.1, 0.2, 0.3},
      xmin=-0.7, xmax=10.7, xtick={0,...,10},
      axis lines*=left,
      xlabel={number of heads $k$ in 10 flips of a fair coin},
      ylabel={probability $P(X = k)$},
      % numbers in the text font (Lato), not the math font
      xticklabel={\pgfmathprintnumber[assume math mode=true]{\tick}},
      tick label style={font=\small}, label style={font=\small},
    ]
    % exact binomial(10, 1/2) probabilities: C(10, k) / 1024
    \addplot[fill=darkblue!65, draw=none] coordinates
      {(0,1/1024) (1,10/1024) (2,45/1024) (3,120/1024) (4,210/1024) (5,252/1024)
       (6,210/1024) (7,120/1024) (8,45/1024) (9,10/1024) (10,1/1024)};
  \end{axis}
\end{tikzpicture}
\end{minipage}%
\hfill
\begin{minipage}[c]{0.41\textwidth}
\centering\small\setlength{\tabcolsep}{3.5pt}
% probabilities rounded from C(10, k) / 1024; the rounded P(X = k) add up to the P(X <= k) column
\begin{tabular}{c|ccc}
  \toprule
  heads $k$ & $\binom{10}{k}$ & $P(X = k)$ & $P(X \le k)$ \\
  \midrule
  0  & 1   & 0.001 & 0.001 \\
  1  & 10  & 0.010 & 0.011 \\
  2  & 45  & 0.044 & 0.055 \\
  3  & 120 & 0.117 & 0.172 \\
  4  & 210 & 0.205 & 0.377 \\
  5  & 252 & 0.246 & 0.623 \\
  6  & 210 & 0.205 & 0.828 \\
  7  & 120 & 0.117 & 0.945 \\
  8  & 45  & 0.044 & 0.989 \\
  9  & 10  & 0.010 & 0.999 \\
  10 & 1   & 0.001 & 1.000 \\
  \bottomrule
\end{tabular}
\end{minipage}

\question{Q1.}
In 10 flips of a fair coin, what is the probability that the result would turn out to be
``[0 or 1 head] or [0 or 1 tail]''?
(i.e., all flips being the same side or with just one exception)
\answerbox[1.2cm]{0 or 1 head, or 0 or 1 tail, means 0, 1, 9, or 10 heads:
$0.001 + 0.010 + 0.010 + 0.001 = 0.022$.}

\medskip
\begin{highlightbox}
A dealer brings a coin to play betting games, and claims that it is a fair coin.
You are suspicious about the coin.
You may watch 10 flips; then you must decide whether to accuse the dealer of using an unfair coin.
\textbf{Before any coin is flipped, plan how you will decide.}
\end{highlightbox}

% ------------------------------------------------------------------ page 2: the plan, then the data
\newpage
\begin{highlightbox}
\textbf{Definitions.}
The \emph{null hypothesis}, $H_0$, is the claim put to the test:
a statement about how the data are generated,
precise enough that the probability of every possible result can be computed.
A \emph{test statistic} is a number computed from the data to make the decision.
Its \emph{sampling distribution under the null hypothesis} is its probability distribution if $H_0$ is true:
how it would vary over many repetitions of the experiment.
\end{highlightbox}

\question{Q2.}
What's your null hypothesis (in English)?
\answerbox[1.1cm]{The coin is fair: each flip lands heads with probability 1/2, independently of the other flips.
(``The coin is unfair'' cannot be $H_0$: it does not fix the probability of any result.)}

\question{Q3.}
You decide to use the number of heads in the 10 flips as the test statistic.
What is its sampling distribution under the null hypothesis?
(Hint: you have been staring at it since page 1. Wink wink.)
\answerbox[0.9cm]{The binomial distribution with $n = 10$ and $p = 1/2$: the chart and table on page 1.}

\question{Q4.}
A \emph{decision rule} lists, before any data are seen,
the values of the test statistic that lead you to accuse the dealer.
Your rule must not wrongly accuse an honest dealer more than 1 time in 20 (a probability of at most 0.05).
Build the rule in three steps.

(a) If the null hypothesis is true, which numbers of heads are least likely? Look at the chart.
\answerbox[0.9cm]{0 and 10 heads (0.001 each), then 1 and 9 (0.010 each), then 2 and 8 (0.044 each).}

(b) Choose the rule that accuses on as many numbers of heads as possible within the limit.
What is the probability that it accuses an honest dealer? Add the probabilities from the table.
\answerbox[1.3cm]{Accuse on 0, 1, 9, or 10 heads: $0.001 + 0.010 + 0.010 + 0.001 = 0.022$.
Adding 2 and 8 makes $0.055 + 0.055 = 0.110$, over the limit.}

(c) Why must the rule be fixed before seeing any flips?
\answerbox[1.2cm]{Otherwise we could always pick a list that contains what we saw,
and the 1-in-20 promise, which is about the rule, would mean nothing.}

\medskip
The dealer flips: \quad \texttt{H H H T H H H H H H} \quad (9 heads, 1 tail).

\question{Q5.}
Apply your decision rule from Q4. Do you accuse?
\answerbox[1.0cm]{Yes: 9 heads is on the list $\{0, 1, 9, 10\}$.}

\question{Q6.}
Q2--Q5 are steps of a general recipe.
Write it as numbered instructions a classmate could follow to test a different claim with different data;
step 1 is done for you. Which steps would they have to redo?
\prefilledbox[3.8cm]{1. State the null hypothesis: the claim to test,
precise enough to compute the probability of every result (here: the coin is fair).}{%
2. Choose a test statistic (the number of heads).\par
3. Find its sampling distribution under the null hypothesis (the binomial table).\par
4. Before seeing the data, list the values that lead to rejecting the null hypothesis,
so that a true null hypothesis is rejected at most a chosen fraction of the time (1 in 20).\par
5. Collect the data, compute the test statistic, and reject if it is on the list.\par
Redo steps 2--4 for a new claim: the statistic, its distribution under the null, and the list.}

% ------------------------------------------------------------------ key only
% Kept off the worksheet pages so that the key's first pages line up with the student copy.
\ifsolutions
\newpage
\section*{Instructor notes}
\small

The student pages define only the null hypothesis, the test statistic,
and its sampling distribution under the null hypothesis (box before Q2).
Everything else the students build in plain words; name it in the debrief:

\smallskip
\begin{center}\small
\begin{tabular}{p{0.5\textwidth}p{0.4\textwidth}}
  \toprule
  In the worksheet & Standard name \\
  \midrule
  ``the coin is unfair'' (board, after Q6) & alternative $H_1$, two-sided or one-sided \\
  wrongly accusing an honest dealer (Q4b) & type I error, false positive \\
  the limit of 1 time in 20 (Q4) & significance level $\alpha = 0.05$ \\
  the decision rule's list of values (Q4b) & rejection region \\
  how often an honest coin is at least as surprising (board) & p-value \\
  accuse when that number is at most the limit (board) & reject $H_0$ when $p \le \alpha$ \\
  the recipe (Q6) & the algorithm of a null hypothesis test \\
  \bottomrule
\end{tabular}
\end{center}

\begin{itemize}[leftmargin=1.2em, itemsep=3pt]
  \item \textbf{Timing.} About 25 minutes of worksheet: Q1--Q4 15, Q5 2, Q6 8;
    then the p-value at the board.
    Protect the time for Q6, and have two students read their recipes aloud;
    in the afternoon, students turn those instructions into Python.
  \item \textbf{Q1} is the same probability that later becomes the false alarm rate (Q4b)
    and the p-value (board, after Q6); do not point it out until then.
  \item \textbf{Q2.} A null hypothesis must fix the probability of every result:
    ``the coin is fair'' does, ``the coin is unfair'' does not, which is why the fair coin is $H_0$.
  \item \textbf{Q4.} The 1-in-20 limit is supplied, not derived.
    It is a property of the rule over many repeated experiments with a fair coin,
    never the probability that a particular accusation is wrong.
    Rules that accuse on a fixed list of head counts cannot hit 5\% exactly; this one runs at 2.1\%.
  \item \textbf{At the board after Q6: discover the p-value.}
    The dealer protests: ``A fair coin can give nine heads!''
    Ask students which results would be at least as surprising as nine heads if the dealer is honest,
    and how often an honest coin gives one of them.
    Both directions gives 0.022 (exactly 22/1024), heads only 0.011 (11/1024):
    this is where two-sided and one-sided tests come from,
    and why the direction must be chosen before the flips.
    Then compare that number with the 1-in-20 limit: the rule accuses exactly when it is at most 0.05.
    Name it now: the p-value. Reject when $p \le \alpha$;
    the p-value is the smallest $\alpha$ at which the observed result would be rejected.
  \item \textbf{Debrief: why a tail, not the observed value alone?}
    Ask why we do not just report the chance of exactly nine heads (0.010).
    In 1000 fair flips, even the most likely single result, exactly 500 heads, has probability 0.0252:
    any single outcome becomes rare as the possible outcomes multiply,
    so surprise is measured by the chance of a result at least as extreme.
  \item \textbf{Debrief: what the number does not mean.}
    It is computed assuming the dealer is honest, so it is not the probability that the dealer is honest.
    Make it concrete: if 1 dealer in 100 cheats, with a coin that lands heads 90\% of the time,
    the rule accuses an honest dealer with probability 0.021 and a cheater with probability 0.736
    (9 or 10 heads, or 0 or 1, at $p = 0.9$).
    Among accused dealers, the fraction that are honest is
    \[
      \frac{0.99 \times 0.021}{0.99 \times 0.021 + 0.01 \times 0.736} \approx 0.74:
    \]
    most accusations are wrong, although the rule is a 5\% rule (Day 1's conditional probability).
    And not accusing does not prove honesty: 8 heads is no accusation, yet an unfair coin easily gives 8 heads.
  \item \textbf{Power, at the board after Q6.} Ask students to predict first:
    if the coin really lands heads with probability 0.7, how often does the rule accuse?
    Almost always, about half the time, or rarely?
    Rarely: $P(X = 9) \approx 0.121$, $P(X = 10) \approx 0.028$, $P(X \le 1) \approx 0.0001$, about 0.15 in total.
    That is the \emph{power} of the test against $p = 0.7$; more flips would raise it,
    which leads to sample-size planning on Day 3.
  \item \textbf{Mouse, in the debrief.} A mouse turns left on 9 of 10 T-maze trials:
    does the same table test ``no side preference''?
    Only if the trials act like independent flips with one probability; mice learn, alternate, and perseverate,
    so step 3 (the sampling distribution under the null) needs a different model.
    Likewise for cells counted from one tube or one animal.
\end{itemize}
\fi

\end{document}
